\section{第一章 C++ 语言入门}
\subsection{1.1 语言概述与基础定位}
\begin{itemize}
\item C++语言诞生与发展历程
\item 面向过程程序设计思想
\item 面向对象程序设计思想
\item C语言兼容特性
\item 高性能运行特点
\item 跨平台可移植特性
\item 程序开发完整流程
\end{itemize}

\subsection{1.2 开发环境与编译链接}
\begin{itemize}
\item 主流C++编译器种类
\item 编译器编译工作原理
\item 集成开发环境功能作用
\item 工程项目目录结构
\item 头文件作用与规范
\item 源文件作用与规范
\item 预处理阶段执行逻辑
\item 编译阶段语法检测
\item 汇编阶段代码转换
\item 链接阶段整合程序
\item 可执行文件生成运行
\end{itemize}

\subsection{1.3 基础语法格式规范}
\begin{itemize}
\item 标识符命名组成规则
\item 标识符合法命名限制
\item 系统关键字定义
\item 系统保留字作用
\item 单行注释书写格式
\item 多行块注释书写格式
\item 程序语句结束标识
\item 代码块边界划分规则
\item 代码缩进排版规范
\end{itemize}

\section{第二章 变量常量与数据类型}
\subsection{2.1 变量完整核心概念}
\begin{itemize}
\item 变量本质存储含义
\item 变量声明语法格式
\item 变量定义内存开辟
\item 变量初始化赋值操作
\item 变量二次赋值数值修改
\item 变量局部作用域范围
\item 变量全局作用域范围
\item 变量内存生命周期
\item 变量销毁释放机制
\end{itemize}

\subsection{2.2 全部基础内置数据类型}
\begin{itemize}
\item 字节短整型 short
\item 基本整型 int
\item 长整型 long
\item 超长整型 long long
\item 单精度浮点型 float
\item 双精度浮点型 double
\item 字符型 char 编码存储
\item 布尔型 bool 逻辑判定
\item 空类型 void 占位作用
\item 各类型取值范围边界
\item 各类型占用字节大小
\end{itemize}

\subsection{2.3 常量分类与定义方式}
\begin{itemize}
\item 字面常量直接数值常量
\item 整型字面常量
\item 浮点字面常量
\item 字符字面常量
\item 字符串字面常量
\item const修饰符号常量
\item 常量不可二次修改特性
\item 常量作用域划分
\end{itemize}

\subsection{2.4 数据类型相互转换规则}
\begin{itemize}
\item 隐式自动类型转换触发条件
\item 低精度向高精度自动提升
\item 强制类型转换语法写法
\item 手动高精度转低精度
\item 类型转换精度丢失现象
\item 转换后数据结果判定
\end{itemize}

\section{第三章 运算符体系全分类}
\subsection{3.1 算术运算符}
\begin{itemize}
\item 加法运算规则
\item 减法运算规则
\item 乘法运算规则
\item 除法整数与小数运算区别
\item 取模余数运算限制条件
\item 前置自增运算执行顺序
\item 后置自增运算执行顺序
\item 前置自减运算执行顺序
\item 后置自减运算执行顺序
\end{itemize}

\subsection{3.2 赋值与复合赋值运算符}
\begin{itemize}
\item 基础赋值运算逻辑
\item 加法复合赋值运算
\item 减法复合赋值运算
\item 乘法复合赋值运算
\item 除法复合赋值运算
\item 取模复合赋值运算
\end{itemize}

\subsection{3.3 关系比较运算符}
\begin{itemize}
\item 等于比较运算
\item 不等于比较运算
\item 大于比较运算
\item 小于比较运算
\item 大于等于比较运算
\item 小于等于比较运算
\item 关系运算布尔结果返回
\end{itemize}

\subsection{3.4 逻辑运算符}
\begin{itemize}
\item 逻辑与运算判定规则
\item 逻辑或运算判定规则
\item 逻辑非运算取反规则
\item 逻辑短路与执行特性
\item 逻辑短路或执行特性
\item 多逻辑组合运算判断
\end{itemize}

\subsection{3.5 位运算符与特殊运算符}
\begin{itemize}
\item 按位与运算规则
\item 按位或运算规则
\item 按位异或运算规则
\item 按位取反运算规则
\item 二进制左移运算
\item 二进制右移运算
\item 三目条件运算符格式
\item 运算符优先级排序规则
\item 运算符左右结合性判定
\end{itemize}

\section{第四章 程序流程控制结构}
\subsection{4.1 顺序执行结构}
\begin{itemize}
\item 自上而下逐行执行逻辑
\item 顺序代码执行顺序规则
\end{itemize}

\subsection{4.2 分支判断结构}
\begin{itemize}
\item if单分支条件判断语句
\item if-else双分支选择语句
\item else if多分支区间匹配
\item switch开关分支语句格式
\item switch常量匹配规则
\item case分支穿透问题处理
\item default默认分支作用
\item 分支语句嵌套使用方式
\end{itemize}

\subsection{4.3 循环遍历结构}
\begin{itemize}
\item while前置条件循环结构
\item while循环执行流程
\item do-while后置循环结构
\item do-while先执行后判断特性
\item for计数循环完整格式
\item for循环初始化条件判断迭代
\end{itemize}

\subsection{4.4 循环跳转与多层嵌套}
\begin{itemize}
\item break语句跳出当前循环
\item break终止多层循环用法
\item continue跳过单次循环迭代
\item 单层循环业务逻辑编写
\item 双层循环嵌套运行规则
\item 多层循环嵌套实际应用
\item 无限循环构成与终止条件
\end{itemize}

\section{第五章 数组与字符串处理}
\subsection{5.1 一维数组完整概念}
\begin{itemize}
\item 数组连续内存存储原理
\item 一维数组定义语法
\item 一维数组初始化方式
\item 数组下标索引从零开始规则
\item 通过下标存取数组元素
\item 数组整体全部遍历读取
\item 数组元素单个修改赋值
\end{itemize}

\subsection{5.2 多维数组结构}
\begin{itemize}
\item 二维数组行列结构含义
\item 二维数组定义与初始化
\item 二维数组下标定位元素
\item 二维数组嵌套遍历方式
\item 多维数组扩展结构特性
\end{itemize}

\subsection{5.3 C++标准字符串操作}
\begin{itemize}
\item string字符串类定义创建
\item 字符串直接赋值初始化
\item 字符串拼接合并操作
\item 字符串截取子串功能
\item 字符串指定字符查找
\item 字符串内容替换修改
\item 字符串分割拆分处理
\item 字符串大小比较判定
\item 字符串长度获取方法
\end{itemize}

\section{第六章 函数编程体系}
\subsection{6.1 函数定义声明与调用}
\begin{itemize}
\item 函数代码封装复用思想
\item 函数声明作用与格式
\item 函数实体定义编写
\item 函数调用触发执行流程
\item 主函数程序入口作用
\item 自定义函数调用规则
\end{itemize}

\subsection{6.2 函数参数传递机制}
\begin{itemize}
\item 形式参数定义特性
\item 实际参数传入规则
\item 形参实参匹配对应
\item 值传递数据拷贝原理
\item 值传递单向数据传递
\item 地址传递内存地址传递
\item 引用传递别名传参特性
\end{itemize}

\subsection{6.3 函数返回值处理}
\begin{itemize}
\item 单个返回值定义格式
\item return语句返回数据
\item 返回值接收存储方式
\item 无返回值函数定义
\item 多数据组合返回处理
\end{itemize}

\subsection{6.4 函数进阶核心特性}
\begin{itemize}
\item 函数重载判定约束条件
\item 同名重载函数区分调用
\item 递归函数自身调用逻辑
\item 递归递进执行过程
\item 递归回溯返回过程
\item 递归终止边界条件设置
\item 递归调用栈内存变化
\end{itemize}

\section{第七章 指针与引用机制}
\subsection{7.1 指针基础底层原理}
\begin{itemize}
\item 程序内存地址编号含义
\item 指针变量存储地址值
\item 取地址运算符使用
\item 解引用取值运算符作用
\item 通过指针修改变量数据
\item 空指针定义与作用
\item 野指针产生原因与风险
\item 指针变量占用内存大小
\end{itemize}

\subsection{7.2 指针组合实操应用}
\begin{itemize}
\item 指针与数组结合访问元素
\item 指针移动遍历数组
\item 指针数组存储地址集合
\item 数组指针指向数组整体
\item 指针作为函数参数传参
\item 多级指针嵌套地址存储
\end{itemize}

\subsection{7.3 引用别名机制详解}
\begin{itemize}
\item 引用变量定义语法格式
\item 引用绑定原有变量规则
\item 引用不可二次重绑定
\item 引用代替指针简化操作
\item 引用传参高效修改实参
\item 引用与指针本质区别
\end{itemize}

\section{第八章 结构体、共用体与枚举}
\subsection{8.1 结构体自定义复合类型}
\begin{itemize}
\item 结构体类型格式定义
\item 结构体内部成员组成
\item 结构体变量创建实例
\item 结构体变量初始化赋值
\item 点运算符访问结构体成员
\item 指针访问结构体成员
\item 结构体数组批量存储对象
\end{itemize}

\subsection{8.2 共用体共享内存类型}
\begin{itemize}
\item 共用体内存共享核心特点
\item 共用体类型定义格式
\item 共用体变量成员赋值
\item 共用体数据覆盖特性
\item 共用体适用业务场景
\end{itemize}

\subsection{8.3 枚举常量限定类型}
\begin{itemize}
\item 枚举类型定义规范
\item 枚举常量默认数值规则
\item 手动修改枚举常量数值
\item 枚举变量取值限定范围
\item 枚举类型实际应用场景
\end{itemize}

\section{第九章 面向对象编程基础}
\subsection{9.1 类与对象核心关系}
\begin{itemize}
\item 类抽象模板概念含义
\item 类内部成员分类组成
\item 对象实例化创建过程
\item 类与对象从属关联关系
\item 成员变量描述对象属性
\item 成员方法描述对象行为
\item 对象调用自身成员规则
\end{itemize}

\subsection{9.2 构造函数与析构函数}
\begin{itemize}
\item 构造函数初始化对象作用
\item 无参构造函数默认规则
\item 有参构造函数自定义初始化
\item 构造函数重载特性
\item 析构函数资源释放功能
\item 对象销毁自动调用析构
\end{itemize}

\subsection{9.3 访问权限与封装特性}
\begin{itemize}
\item public公有访问权限范围
\item private私有访问权限范围
\item protected保护访问权限范围
\item 类内部成员访问规则
\item 类外部成员访问限制
\item 数据封装隐藏设计思想
\item 公有接口私有数据规范
\end{itemize}

\section{第十章 面向对象编程进阶}
\subsection{10.1 类的继承机制}
\begin{itemize}
\item 父类基类定义概念
\item 子类派生类定义概念
\item 单继承父子类从属关系
\item 继承后成员自动拥有规则
\item 不同权限成员继承访问
\item 子类构造函数调用父类构造
\item 继承代码复用优势特点
\end{itemize}

\subsection{10.2 成员方法重写}
\begin{itemize}
\item 子类重写父类方法条件
\item 重写方法语法格式约束
\item 重写后覆盖原有父类功能
\item 子类自身方法与重写区分
\end{itemize}

\subsection{10.3 多态实现原理}
\begin{itemize}
\item 虚函数定义语法格式
\item 虚函数表底层存储机制
\item 对象向上转型形态转换
\item 父类指针指向子类对象
\item 多态动态绑定调用方法
\item 多态代码扩展适配特性
\end{itemize}

\section{第十一章 内存分区与动态内存管理}
\subsection{11.1 程序内存五大分区}
\begin{itemize}
\item 栈区临时变量存储特性
\item 栈内存自动分配释放
\item 堆区动态内存空间特点
\item 全局区全局变量存储
\item 静态区静态变量生命周期
\item 常量区只读数据存储规则
\item 代码区程序指令存储
\end{itemize}

\subsection{11.2 动态内存操作关键字}
\begin{itemize}
\item new关键字开辟堆内存
\item new创建单个动态变量
\item new创建动态数组空间
\item delete关键字释放内存
\item delete释放单个堆空间
\item delete[]释放动态数组
\item 动态内存泄漏产生原因
\item 动态内存合理使用规范
\end{itemize}

\section{第十二章 模板通用编程}
\subsection{12.1 函数模板}
\begin{itemize}
\item 模板参数通用类型定义
\item 函数模板完整定义格式
\item 模板函数自动类型推导
\item 函数模板适配多种数据
\item 模板函数调用执行规则
\end{itemize}

\subsection{12.2 类模板}
\begin{itemize}
\item 类模板通用结构封装
\item 类模板成员定义写法
\item 模板类实例化创建对象
\item 模板类多参数定义用法
\end{itemize}

\section{第十三章 STL标准模板库}
\subsection{13.1 STL整体组成架构}
\begin{itemize}
\item STL标准库设计优势
\item 容器存储数据作用
\item 迭代器定位遍历作用
\item 通用算法批量处理数据
\end{itemize}

\subsection{13.2 序列式常用容器}
\begin{itemize}
\item vector向量动态数组特性
\item vector增删查改基础操作
\item list双向链表容器结构
\item list链表元素存取特点
\end{itemize}

\subsection{13.3 关联式常用容器}
\begin{itemize}
\item set集合自动去重原理
\item set有序排序存储规则
\item map键值成对存储结构
\item map键唯一值可重复特性
\end{itemize}

\subsection{13.4 迭代器遍历机制}
\begin{itemize}
\item 迭代器本质模拟指针
\item 迭代器定义与指向元素
\item 正向迭代器遍历容器
\item 常量迭代器只读限制
\end{itemize}

\section{第十四章 文件读写操作}
\subsection{14.1 文件流基础分类与操作}
\begin{itemize}
\item 文件输入流读取数据
\item 文件输出流写入数据
\item 文件打开语法与模式选择
\item 文件正常关闭规范操作
\item 文件路径绝对相对写法
\end{itemize}

\subsection{14.2 文本文件与二进制文件}
\begin{itemize}
\item 文本文件字符读写方式
\item 二进制字节原始数据读写
\item 单行多行文本读取处理
\item 批量数据写入保存文件
\end{itemize}

\section{第十五章 异常处理机制}
\subsection{15.1 异常抛出与捕获处理}
\begin{itemize}
\item 程序异常产生触发场景
\item throw语句主动抛出异常
\item try监控异常代码块
\item catch匹配对应异常类型
\item 异常捕获后业务处理
\end{itemize}

\subsection{15.2 异常匹配与层级传递}
\begin{itemize}
\item 多异常类型分级捕获顺序
\item 通用异常统一捕获处理
\item 函数内部异常向上传递
\item 多层调用异常逐层抛出
\end{itemize}
